\chapter[One Evidence Bound and Local-Collective Free Energy]{One Evidence Bound and\\Local-Collective Free Energy}
\label{ch:elbo}
\label{ch:local-collective-elbo}

\Cref{ch:generative} fixed one normalized joint law and \cref{ch:probability}
one correlated recognition coupling of the whole population. This chapter asks
what exact relation binds the two, and how the objective of one agent, or of a
block of agents, is related to that of the population. The answer is one
extended evidence bound, defined for every probability law on the latent space
(\cref{def:elbo-extended,thm:elbo-extended-gap}), and the local--global
potential identity (\cref{thm:obs-local-global-potential}): with the outside
marginal held fixed, the change of the conditional local free energy of a block
equals the change of the collective free energy, so the agents' local
objectives form an exact potential game (\cref{cor:obs-potential-game}). Every
local object is a disintegration of the one bound through
\cref{lem:prob-kl-chain-rule}. The bound is the objective of
\cref{ch:selective-communication,ch:relational-inference,ch:two-timescales},
and \cref{ch:coarsegraining,ch:full-graph-meta-agent-vfe} transport it along
evidence-preserving channels.

\paragraph{Scope.}
The chapter does not claim that a maximizer of the bound over a restricted
recognition family exists, that the bound is tight for any particular
restricted family, that the alternating scheme of exact coordinates converges
(\cref{open:elbo-alternating-convergence}), or that a block of agents together
with its local free energy is a meta-agent. No criterion internal to one bound
that selects a recognition structure is known
(\cref{open:restrict-intrinsic-criterion}). \status{NOT-CLAIMED}

\section{Recognition coupling and total correlation}
\label{sec:elbo-recognition}

The population recognition law is the coupling
$\mathbb Q_{V,o,X}\in\operatorname{Cpl}(\mathbb Q_{\bullet,o,X})$ selected by
\cref{hyp:prob-measurable-coupling-selector}, and every restriction used later
is a subfamily $\mathcal Q'$ of \cref{prop:restrict-principle}. Its marginals
do not determine it once two of them are non-Dirac
(\cref{prop:prob-marginals-do-not-determine-joint}), so they cannot replace it
in an exact entropy or evidence calculation; the next two statements give the
error its exact value and sign.

\corollaryheading{Extended total-correlation chain identity}{thm:elbo-total-correlation-chain}
Let $Y=(Y_j)_{j\in\mathcal J}$ take values in a finite product of standard
Borel spaces, let $Q$ be a probability law with marginals $Q_j$, and let
$(\varrho_j)_{j\in\mathcal J}$ be arbitrary probability laws. With the total
correlation
\begin{equation}
\operatorname{TC}(Q):=\KL\Bigl(Q\Bigm\Vert\bigotimes_{j}Q_j\Bigr)\in[0,+\infty],
\label{eq:obs-total-correlation}
\end{equation}
the identity
\begin{equation}
\KL\Bigl(Q\Bigm\Vert\bigotimes_j\varrho_j\Bigr)
=\operatorname{TC}(Q)+\sum_j\KL(Q_j\Vert\varrho_j)
\label{eq:elbo-total-correlation-chain}
\end{equation}
holds in $[0,+\infty]$, including every singular branch. The infinite branch
occurs for a continuous diagonal: if $U\sim\operatorname{Unif}[0,1]$ and
$Q=\operatorname{law}(U,U)$, both marginals are uniform and
$\operatorname{TC}(Q)=+\infty$. \status{ESTABLISHED}

\noindent\emph{Proof.}
The identity is \eqref{eq:prob-kl-product-reference} of
\cref{lem:prob-kl-chain-rule} with its first term named. The diagonal is null
for the uniform law on the square and carries all the mass of $Q$, so $Q$ is
singular with respect to $Q_1\otimes Q_2$. \hfill$\square$

\propositionheading{Finite-density total correlation and the two signs}{prop:elbo-total-correlation-signs}
Partition the latent coordinates into blocks $j$, let $\mathbb Q_{V,o,X}$
have density $q_{V,o,X}$ with respect to $\nu^{\mathsf Y}$ and block marginal
densities $q_j$, and suppose the joint and block differential entropies are
finite. Then the total correlation
$\operatorname{TC}(\mathbb Q_{V,o,X})\geq0$ of the blocks,
\eqref{eq:obs-total-correlation}, satisfies
\begin{equation}
\E_{\mathbb Q_{V,o,X}}\bigl[\log q_{V,o,X}(Y)\bigr]
=\sum_j\E_{\mathbb Q_{V,o,X}}\bigl[\log q_j(Y_j)\bigr]
+\operatorname{TC}(\mathbb Q_{V,o,X}),
\label{eq:elbo-entropy-defect}
\end{equation}
with equality of the two log-density expectations exactly when the blocks are
independent. If in addition
$\E_{\mathbb Q_{V,o,X}}\lvert\log p_{\theta,V}(o,Y\given X)\rvert<\infty$, the
substituted free energy and its pseudo-ELBO,
\begin{equation}
\widetilde\Fenergy
:=\E_{\mathbb Q_{V,o,X}}\Bigl[-\log p_{\theta,V}(o,Y\given X)
+\sum_j\log q_j(Y_j)\Bigr],
\qquad
\widetilde\Lelbo:=-\widetilde\Fenergy,
\label{eq:elbo-substituted-free-energy}
\end{equation}
satisfy
\begin{equation}
\widetilde\Fenergy=\Fenergy_V-\operatorname{TC}(\mathbb Q_{V,o,X}),
\qquad
\widetilde\Lelbo=\Lelbo_V(\mathbb Q_{V,o,X};X,o)+\operatorname{TC}(\mathbb Q_{V,o,X}),
\label{eq:elbo-total-correlation-signs}
\end{equation}
where $\Lelbo_V$ and $\Fenergy_V$ are the classical functionals of
\cref{thm:elbo-exact-identity}. The pseudo-ELBO exceeds the ELBO by the total
correlation and need not bound the log evidence. \status{ESTABLISHED}

\noindent\emph{Proof.}
Take $\varrho_j=\mathbb Q_{V,o,X,j}$ in
\cref{thm:elbo-total-correlation-chain}; the entropy hypotheses make every
term finite and allow the logarithm of $q_{V,o,X}/\prod_jq_j$ to be split,
which is \eqref{eq:elbo-entropy-defect}. Subtracting
\eqref{eq:elbo-substituted-free-energy} from
$\Fenergy_V=\E_{\mathbb Q_{V,o,X}}[-\log p_{\theta,V}+\log q_{V,o,X}]$ leaves
$\operatorname{TC}$, and negation gives the second identity. \hfill$\square$

\section{Evidence domain}
\label{sec:elbo-hypotheses}

\hypothesisheading{Evidence-domain hypotheses}{hyp:elbo-evidence-domain}
\label{hyp:prob-regular-observation}
\label{hyp:prob-classical-split-elbo}
The extended bound of \cref{sec:elbo-identity} uses (H1)--(H2); its classical
split into two expected log densities also uses (H3)--(H4).

\textbf{(H1) Types and version.} The joint
$\mathbb P_{\theta,V}(\cdot\given X)$ of \cref{ch:generative} has the density
$p_{\theta,V}(o,y\given X)$ relative to its reference, and
$\boldsymbol\Pi_{\theta,V,o,X}$ is the version \eqref{eq:prob-rcp-density} on
the regular set $\mathsf O^{\mathrm{reg}}_{\theta,V,X}$ of
\cref{hyp:prob-regular-conditional-convention}, so every posterior and bound
below holds for $\mathbb P^{\mathsf O}_{\theta,V,X}$-almost every $o$. A
comparison of finitely many parameter values uses an observation in the
intersection of their regular sets. No density of the recognition law is
required.

\textbf{(H2) Positive finite evidence.}
$0<p_{\theta,V}(o\given X)<\infty$; this excludes the observations at which
the posterior formula would divide by zero or by infinity.

\textbf{(H3) Absolute continuity.}
$\mathbb Q_{V,o,X}\ll\boldsymbol\Pi_{\theta,V,o,X}$. Since the posterior is
dominated by $\nu^{\mathsf Y}$, the recognition law then has a density
$q_{V,o,X}=d\mathbb Q_{V,o,X}/d\nu^{\mathsf Y}$, with a jointly measurable
version for parameterized kernels under
\cref{thm:prob-kernel-rn-measurable-version,cor:prob-common-reference-joint-density}.
Clause (H3) excludes every recognition law that charges a posterior-null set;
such a law lies on the singular branch of the extended bound, where the bound
equals $-\infty$ (\cref{prop:elbo-subspace-support-singular}).

\textbf{(H4) Log-integrability.} The generative and recognition halves
\begin{align}
&\E_{\mathbb Q_{V,o,X}}\bigl\lvert\log p_{\theta,V}(o,Y\given X)\bigr\rvert<\infty,
\label{eq:elbo-log-integrability}\\
&\E_{\mathbb Q_{V,o,X}}\bigl\lvert\log q_{V,o,X}(Y)\bigr\rvert<\infty,
\label{eq:prob-integrability}
\end{align}
both hold. Clause (H4) excludes recognition laws whose tails make the two
classical expectations individually infinite. It is not implied by
(H1)--(H3), and the extended bound is well defined without it;
\cref{prop:elbo-integrability-witnesses} gives a Cauchy witness whose formal
difference is $-\infty$ and a countable witness whose two expectations are
both $-\infty$.

\propositionheading{Posterior-null support has infinite cost}{prop:elbo-subspace-support-singular}
Let $Q$ be a probability law on the latent space with $Q(S)=1$ and
$\boldsymbol\Pi_{\theta,V,o,X}(S)=0$ for some measurable $S$. Then (H3) fails
for $Q$, $\KL(Q\Vert\boldsymbol\Pi_{\theta,V,o,X})=+\infty$, and
$\Lelbo^{\mathrm{ext}}_V(Q;X,o)=-\infty$. In particular this holds when the
posterior is dominated by Lebesgue measure on $\R^d$ and $Q$ is supported on a
proper affine subspace. A proper closed support alone does not suffice: the
uniform law on $[0,\tfrac12]$ has divergence $\log2$ from the uniform law on
$[0,1]$. \status{ESTABLISHED}

\noindent\emph{Proof.}
The first claims are the null-set clause of
\cref{prop:prob-density-absolute-continuity} and
\eqref{eq:elbo-extended-definition}. A proper affine subspace is Lebesgue
null, and the uniform law on $[0,\tfrac12]$ has density $2$ on its support
relative to the uniform law on $[0,1]$. \hfill$\square$

\section{The extended evidence bound and its record presentation}
\label{sec:elbo-identity}
\label{sec:local-interaction-record-model}
\label{sec:local-collective-vfe}

\hypothesisheading{Agent state and interaction records}{hyp:local-interaction-kernels}
Agent $i$ has a block space $(\mathsf Y_i,\mathscr Y_i)$, and
$(\mathsf Y,\mathscr Y)=\prod_{i\in V}(\mathsf Y_i,\mathscr Y_i)$. The block may
contain auxiliary latent variables and law points of $(\mathcal E_b)_{c_*}$ and
$(\mathcal E_m)_{c_*}$, which are then latent coordinates at the law-point or
configuration level of \cref{def:prob-description-levels}, not the recognition
marginal of agent $i$. The baseline is the latent population kernel
$\mathbb P^{\mathsf Y}_{\theta,V}:(\mathsf X,\mathscr X)\rightsquigarrow(\mathsf Y,\mathscr Y)$
of \cref{cons:gen-finite-directed-law}. Each record $a\in\mathfrak A$ occurs
once and has a finite scope $\partial a\subseteq V$, a record space
$(\mathsf O_a,\mathscr O_a)$, and a record kernel $K_a$ of the type
\eqref{eq:prob-record-kernel-signature}, all of whose values
$K_a(X,y_{\partial a},\cdot)$ are dominated by one $\sigma$-finite measure
$\nu_a$ with a jointly measurable negative log density $E_{a,o}$ in
$(-\infty,+\infty]$,
\begin{equation}
K_a(X,y_{\partial a},do_a)=e^{-E_{a,o}(X,y_{\partial a})}\nu_a(do_a),
\label{eq:obs-interaction-density}
\end{equation}
where $E_{a,o}$ depends on $o$ only through $o_a$. A record shared by agents
$i$ and $j$ is one element of $\mathfrak A$, and writing both $E_{ij}$ and
$E_{ji}$ asserts two directed records. Binary success records
$K_a(y_{\partial a},\{1\})=e^{-E(y_{\partial a})}$ with measurable
$E:\mathsf Y_{\partial a}\to[0,+\infty]$ are admitted. The kernels and
$\mathbb P^{\mathsf Y}_{\theta,V}$ read neither a recognition law, a
variational parameter, nor a posterior (\cref{req:gen-typing-prohibition}),
and under a common vertex gauge the coordinate presentation of each $K_a$ is
equivariant under the simultaneous change of the incident frames, as in
\eqref{eq:gen-gauge-record-kernel}, so the kernel itself is frame independent.
The hypothesis excludes record families without a common dominating measure
(\cref{open:prob-nondominated-likelihood}) and every generative factor that
reads recognition.

Correlated baseline states are permitted; a graph-local Markov blanket needs in
addition the corresponding locality of the regular conditionals of
$\mathbb P^{\mathsf Y}_{\theta,V}$. Frame independence can be realized by giving
record $a$ its own frame and a link $\Theta^x_{a\leftarrow i}$ on each leg, so
that the density reads $\rho_x(\Theta^x_{a\leftarrow i})y^x_i$; comparing the
belief and model channels still requires the cross morphisms of
\cref{sec:geo-cross-morphisms}. From here on one $X\in\mathsf X$ is fixed and
suppressed from $K_a$, $E_{a,o}$ and $\mathbb P^{\mathsf Y}_{\theta,V}$.

\definitionheading{Record likelihood, slice measure, and extended bound}{def:elbo-extended}
\label{thm:obs-collective-vfe}
Assume \cref{hyp:local-interaction-kernels} and (H1)--(H2), and fix
$o\in\mathsf O^{\mathrm{reg}}_{\theta,V,X}$. The record likelihood and the
interaction energy are
\begin{align}
L_o(y)&:=\prod_{a\in\mathfrak A}e^{-E_{a,o}(y_{\partial a})}\in[0,+\infty),
\label{eq:obs-complete-record-likelihood}\\
H_o(y)&:=-\log L_o(y)=\sum_{a\in\mathfrak A}E_{a,o}(y_{\partial a})
\in(-\infty,+\infty];
\label{eq:obs-interaction-energy}
\end{align}
$H_o$ has one term for each record, never equals $-\infty$ because each factor
density is finite, and can be negative. With $\nu^{\mathsf O}=\bigotimes_a\nu_a$
the joint, normalized on every finite record hypergraph, cyclic or not, by
\cref{prop:gen-exact-normalization}, is
\begin{equation}
\mathbb P_{\theta,V}(dy,do\given X)
=\mathbb P^{\mathsf Y}_{\theta,V}(dy)\prod_{a\in\mathfrak A}K_a(y_{\partial a},do_a)
=\mathbb P^{\mathsf Y}_{\theta,V}(dy)L_o(y)\nu^{\mathsf O}(do).
\label{eq:obs-interaction-joint}
\end{equation}
The evidence and the posterior slice are
\begin{align}
p_{\theta,V}(o\given X)&=\int_{\mathsf Y}L_o(y)\mathbb P^{\mathsf Y}_{\theta,V}(dy)\in(0,+\infty),
\label{eq:obs-global-evidence}\\
\boldsymbol\Pi_{\theta,V,o,X}(dy)&=p_{\theta,V}(o\given X)^{-1}L_o(y)\mathbb P^{\mathsf Y}_{\theta,V}(dy),
\label{eq:obs-global-posterior}
\end{align}
and the slice measure is
\begin{equation}
M_{\theta,V,o,X}(S):=\int_Sp_{\theta,V}(o,y\given X)\nu^{\mathsf Y}(dy)
=\int_SL_o(y)\mathbb P^{\mathsf Y}_{\theta,V}(dy)
=p_{\theta,V}(o\given X)\boldsymbol\Pi_{\theta,V,o,X}(S).
\label{eq:elbo-slice-measure}
\end{equation}
For every probability law $Q$ on $(\mathsf Y,\mathscr Y)$ the extended ELBO and
free energy, which are well-defined extended reals because
$\log p_{\theta,V}(o\given X)$ is finite and relative entropy lies in
$[0,+\infty]$, are
\begin{align}
\Lelbo_V^{\mathrm{ext}}(Q;X,o)&:=\log p_{\theta,V}(o\given X)
-\KL\bigl(Q\Vert\boldsymbol\Pi_{\theta,V,o,X}\bigr)
\in\bigl[-\infty,\log p_{\theta,V}(o\given X)\bigr],
\label{eq:elbo-extended-definition}\\
\Fenergy_V^{\mathrm{ext}}[Q;X,o]&:=-\Lelbo_V^{\mathrm{ext}}(Q;X,o)
=-\log p_{\theta,V}(o\given X)+\KL\bigl(Q\Vert\boldsymbol\Pi_{\theta,V,o,X}\bigr),
\label{eq:elbo-extended-free-energy}
\intertext{and with $\theta$, $V$ and $X$ fixed the collective free energy is written}
\Fenergy^{\mathrm{ext}}_o(Q)&:=-\Lelbo^{\mathrm{ext}}_o(Q):=\Fenergy_V^{\mathrm{ext}}[Q;X,o].
\label{eq:obs-collective-vfe}
\end{align}
No domination, density, or mean-field condition on $Q$ enters.

\corollaryheading{Extended gap identity and energy form}{thm:elbo-extended-gap}
For every probability law $Q$ on the latent space,
\begin{equation}
\log p_{\theta,V}(o\given X)-\Lelbo_V^{\mathrm{ext}}(Q;X,o)
=\Fenergy_V^{\mathrm{ext}}[Q;X,o]+\log p_{\theta,V}(o\given X)
=\KL\bigl(Q\Vert\boldsymbol\Pi_{\theta,V,o,X}\bigr),
\label{eq:elbo-extended-gap}
\end{equation}
so $\Lelbo_V^{\mathrm{ext}}\leq\log p_{\theta,V}(o\given X)$ with equality
exactly when $Q=\boldsymbol\Pi_{\theta,V,o,X}$ as measures. If moreover
$Q\ll\mathbb P^{\mathsf Y}_{\theta,V}$,
$\KL(Q\Vert\mathbb P^{\mathsf Y}_{\theta,V})<\infty$, and the negative part of
$H_o$ is $Q$-integrable, then
\begin{equation}
\Fenergy^{\mathrm{ext}}_o(Q)=\KL\bigl(Q\Vert\mathbb P^{\mathsf Y}_{\theta,V}\bigr)+\E_QH_o
\quad\text{in }(-\infty,+\infty].
\label{eq:obs-collective-energy-form}
\end{equation}
No formula $\log p_{\theta,V}=\Lelbo_V^{\mathrm{ext}}+\KL$ is asserted on the
singular branch, where it would read $-\infty+\infty$. \status{ESTABLISHED}

\noindent\emph{Proof.}
The gap identity adds or subtracts the finite number
$\log p_{\theta,V}(o\given X)$ in
\eqref{eq:elbo-extended-definition}--\eqref{eq:elbo-extended-free-energy}, and
the equality case is the zero case of relative entropy. For the energy form
let $r'=dQ/d\mathbb P^{\mathsf Y}_{\theta,V}$. If $Q\ll M_{\theta,V,o,X}$, then
$dQ/dM_{\theta,V,o,X}=r'e^{H_o}$ $Q$-almost surely, and
\cref{prop:elbo-relative-log-representation} gives
$\Fenergy^{\mathrm{ext}}_o(Q)=\int(\log r'+H_o)dQ$; the first term integrates
to the finite $\KL(Q\Vert\mathbb P^{\mathsf Y}_{\theta,V})$ and the second is
quasi-integrable, so the integral splits. If $Q\not\ll M_{\theta,V,o,X}$, some
set null for $M_{\theta,V,o,X}$ has positive $Q$-mass; since
$M_{\theta,V,o,X}$ and $\mathbb P^{\mathsf Y}_{\theta,V}$ agree in null sets on
$\{L_o>0\}$ and $Q\ll\mathbb P^{\mathsf Y}_{\theta,V}$, that mass lies in
$\{H_o=+\infty\}$, and both sides equal $+\infty$. \hfill$\square$

\theoremheading{Classical finite evidence identity}{thm:elbo-exact-identity}
\label{cor:elbo-bound-tightness}
Under \cref{hyp:elbo-evidence-domain}, the extended bound has the separately
integrable representation
\begin{align}
\Lelbo_V(\mathbb Q_{V,o,X};X,o)
&:=\E_{\mathbb Q_{V,o,X}}\bigl[\log p_{\theta,V}(o,Y\given X)-\log q_{V,o,X}(Y)\bigr]
=\Lelbo_V^{\mathrm{ext}}(\mathbb Q_{V,o,X};X,o),
\label{eq:elbo-definition}\\
\Fenergy_V[\mathbb Q_{V,o,X};X,o]&:=-\Lelbo_V(\mathbb Q_{V,o,X};X,o),
\label{eq:elbo-free-energy}
\end{align}
and for $\mathbb P^{\mathsf O}_{\theta,V,X}$-almost every observation
\begin{equation}
\log p_{\theta,V}(o\given X)=\Lelbo_V(\mathbb Q_{V,o,X};X,o)
+\KL\bigl(\mathbb Q_{V,o,X}\Vert\boldsymbol\Pi_{\theta,V,o,X}\bigr).
\label{eq:elbo-identity}
\end{equation}
Consequently $\Lelbo_V\leq\log p_{\theta,V}(o\given X)$, with equality if and
only if $\mathbb Q_{V,o,X}=\boldsymbol\Pi_{\theta,V,o,X}$ as probability
measures on $(\mathsf Y,\mathscr Y)$. Equality of all coordinate marginals is
not sufficient. \status{ESTABLISHED}

\noindent\emph{Proof.}
The posterior and the slice measure have the same null sets, so (H3) gives
$\mathbb Q_{V,o,X}\ll M_{\theta,V,o,X}$, and since the two measures have the
densities $q_{V,o,X}$ and $p_{\theta,V}(o,\cdot\given X)$ relative to
$\nu^{\mathsf Y}$, $d\mathbb Q_{V,o,X}/dM_{\theta,V,o,X}=q_{V,o,X}/p_{\theta,V}(o,\cdot\given X)$
$\mathbb Q_{V,o,X}$-almost surely.
\Cref{prop:elbo-relative-log-representation} gives
$\Lelbo_V^{\mathrm{ext}}=-\int\log(q_{V,o,X}/p_{\theta,V}(o,\cdot\given X))d\mathbb Q_{V,o,X}$,
and (H4) splits the integral, which is \eqref{eq:elbo-definition}. The bound is
then finite, so \eqref{eq:elbo-extended-definition} rearranges to
\eqref{eq:elbo-identity}, and the equality case is that of
\cref{thm:elbo-extended-gap}. Distinct measures with equal coordinate marginals
are supplied by \cref{prop:prob-marginals-do-not-determine-joint}.
\hfill$\square$

\section{Restriction principle}
\label{sec:restrict-principle}

\propositionheading{Restriction principle}{prop:restrict-principle}
Fix $(\theta,X)$ and $o$ as in \cref{def:elbo-extended}, and write
$P^\star=\boldsymbol\Pi_{\theta,V,o,X}$, so that for every law $Q$
\begin{equation}
\Lelbo_V^{\mathrm{ext}}(Q;X,o)=\log p_{\theta,V}(o\given X)-\KL(Q\Vert P^\star).
\label{eq:restrict-exact-identity}
\end{equation}
For arbitrary sets of probability laws $\mathcal Q'\subseteq\mathcal Q$,
\begin{equation}
\sup_{Q\in\mathcal Q'}\Lelbo_V^{\mathrm{ext}}(Q;X,o)
=\log p_{\theta,V}(o\given X)-\inf_{Q\in\mathcal Q'}\KL(Q\Vert P^\star),
\label{eq:restrict-sup-identity}
\end{equation}
the supremum over $\mathcal Q'$ is at most the supremum over $\mathcal Q$, and
when both infima are finite
\begin{equation}
\sup_{\mathcal Q}\Lelbo_V^{\mathrm{ext}}-\sup_{\mathcal Q'}\Lelbo_V^{\mathrm{ext}}
=\inf_{\mathcal Q'}\KL(Q\Vert P^\star)-\inf_{\mathcal Q}\KL(Q\Vert P^\star)\geq0.
\label{eq:restrict-loss}
\end{equation}
No attainment of the suprema or infima is required. \status{ESTABLISHED}

\noindent\emph{Proof.}
Identity \eqref{eq:restrict-exact-identity} is
\eqref{eq:elbo-extended-definition}. Taking suprema turns a finite constant
minus a nonnegative extended divergence into that constant minus its infimum,
including the value $-\infty$. Inclusion gives monotonicity, and subtraction is
legitimate only when both infima are finite. \hfill$\square$

The evidence in \eqref{eq:restrict-exact-identity} is constant because the
generative kernel does not read $Q$ (\cref{req:gen-typing-prohibition}); the
next two propositions show what fails when a model moves with $Q$.

\remarkheading{Orientation, separation, and attainment}{rem:restrict-attainment}
The infimum in \eqref{eq:restrict-sup-identity} is a reverse information
projection, since $Q$ is the first argument of the divergence. Excluding
$P^\star$ from a family does not make the optimum positive or attained: for
$P^\star=\mathcal N(0,1)$ the family $\{\mathcal N(\mu,1):\mu>0\}$ has gaps
$\mu^2/2$ with unattained infimum zero. A positive optimized gap is the
condition $\inf_{Q\in\mathcal Q}\KL(Q\Vert P^\star)>0$. A family constrained to
product laws need not be convex, because mixtures of product laws need not be
product laws. Strict convexity of relative entropy on the ambient law space
therefore does not establish uniqueness on that family; on a convex feasible
family it does give at most one finite minimizer. \status{ESTABLISHED}

\propositionheading{No distinguished target}{prop:gen-no-distinguished-target}
Let $\mathcal Q$ be a nonempty family of recognition laws, and replace one
factor of \eqref{eq:gen-composition} by a nonconstant $Q$-indexed family, giving
a normalized joint $P_{\theta,Q}$ for each $Q\in\mathcal Q$. Assume at the
observation that $0<p_{\theta,Q}(o\given X)<\infty$, that
$Q\ll P_{\theta,Q}(\cdot\given o,X)$, and that the log densities defining the
ELBO are integrable. Then for each fixed $Q$
\begin{equation}
\log p_{\theta,Q}(o\given X)=\Lelbo(Q;X)
+\KL\bigl(Q\Vert P_{\theta,Q}(\cdot\given o,X)\bigr),
\label{eq:gen-moving-target}
\end{equation}
and with $\mathsf e(Q)=\log p_{\theta,Q}(o\given X)$ two members satisfy
\begin{equation}
\Lelbo(Q';X)-\Lelbo(Q;X)
=\mathsf e(Q')-\mathsf e(Q)
-\bigl[\KL\bigl(Q'\Vert P_{\theta,Q'}(\cdot\given o,X)\bigr)
-\KL\bigl(Q\Vert P_{\theta,Q}(\cdot\given o,X)\bigr)\bigr].
\label{eq:gen-elbo-difference}
\end{equation}
On a level set of $\mathsf e$ the change of the bound is the negative change of
the divergence; off it the two changes differ by $\mathsf e(Q')-\mathsf e(Q)$.
The family selects no joint as the model. \status{ESTABLISHED}

\noindent\emph{Proof.}
Apply \cref{thm:elbo-exact-identity} to each fixed pair $(Q,P_{\theta,Q})$ and
subtract. No member is distinguished because no selection rule was declared.
\hfill$\square$

\propositionheading{The stronger consequences fail}{prop:gen-moving-target-witness}
Nonconstant dependence of $P_{\theta,Q}$ on $Q$ does not imply that $\mathsf e$
varies, that no ascent improves a fixed quantity, or that no single member's
evidence bounds the family of bounds. Explicitly, for one binary latent $y$, one
binary observation, and $Q_{\mathrm p}=\operatorname{Ber}(\mathrm p)$ with
$\mathrm p\in(0,1)$, the joint
\begin{equation}
P_{\theta,Q_{\mathrm p}}(o,y\given X)=\tfrac12\operatorname{Ber}\bigl(\tfrac12\mathrm p+\tfrac14\bigr)(y)
\label{eq:gen-moving-target-witness}
\end{equation}
is normalized and injective in $\mathrm p$, has $\mathsf e(Q_{\mathrm p})=\log\tfrac12$
for every $\mathrm p$, and has gaps
\begin{equation}
\tfrac9{10}\log\tfrac97+\tfrac1{10}\log\tfrac13=0.1163218\ldots,
\qquad
\tfrac7{10}\log\tfrac76+\tfrac3{10}\log\tfrac34=0.0216009\ldots,
\qquad
0
\label{eq:gen-moving-target-values}
\end{equation}
at $\mathrm p=\tfrac9{10},\tfrac7{10},\tfrac12$. \status{ESTABLISHED}

\noindent\emph{Proof.}
The observation marginal is $\tfrac12$ for every $\mathrm p$, and the posterior
is $\operatorname{Ber}(\tfrac12\mathrm p+\tfrac14)$, so the gaps are the
Bernoulli relative entropies displayed. The bound increases to the fixed
evidence along the three values, and the member at $\mathrm p=\tfrac12$ has
evidence equal to the supremum of the bounds. \hfill$\square$

\propositionheading{ELBO values do not order unrelated joints}{prop:restrict-unrelated-joint-elbo}
Let $P_\theta$ and $P_{\theta'}$ be normalized joints over the same observation,
possibly with different latent inventories, with recognition families
$\mathcal Q$ and $\mathcal Q'$ on their latent spaces, and let $o$ lie in the
regular set of each. Suppose both infima of the posterior relative entropy
over the respective recognition families are finite. If no evidence-preserving
channel relates the two joints, their optimized bounds satisfy
\begin{equation}
\Lelbo^\star-\Lelbo'^\star
=\bigl[\log p_\theta(o\given X)-\log p_{\theta'}(o\given X)\bigr]
-\Bigl[\inf_{\mathcal Q}\KL-\inf_{\mathcal Q'}\KL\Bigr],
\label{eq:restrict-cross-model}
\end{equation}
which combines an evidence difference, a model-comparison quantity, with a
tightness difference, a recognition-quality quantity. The sign of the left side
orders neither the models nor the recognition families: with binary $O$ and $Y$,
constants $\mathrm c_1\leq0\leq\mathrm c_2$, and the normalized joints
$P(o,1)=e^{\mathrm c_1-\mathrm c_2}$, $P(o,0)=e^{\mathrm c_1}(1-e^{-\mathrm c_2})$,
$P(\bar o,0)=1-e^{\mathrm c_1}$, $P(\bar o,1)=0$ under the singleton family
$\{\delta_1\}$, the pairs $(\mathrm c_1,\mathrm c_2)=(-1,0),(0,2)$ give bounds
$-1>-2$ with the reverse evidence order, and the pairs $(0,1),(-2,0)$ give bounds
$-1>-2$ with the reverse tightness order. \status{ESTABLISHED}

\noindent\emph{Proof.}
Apply \eqref{eq:restrict-sup-identity} to each model and subtract. In the
witness, $\log P(O=o)=\mathrm c_1$ and $P(Y=1\given o)=e^{-\mathrm c_2}$, so the
minimized gap of $\{\delta_1\}$ is $\mathrm c_2$ and the optimized bound is
$\mathrm c_1-\mathrm c_2$. An unused deterministic coordinate appended to either
model, with its point mass appended to the recognition law, makes the latent
inventories literally different and changes no evidence, divergence, or bound.
\hfill$\square$

\openproblemheading{An intrinsic criterion}{open:restrict-intrinsic-criterion}
No criterion internal to one evidence bound at a fixed model is known that
selects a nondegenerate recognition structure. Along nested recognition
families the optimized bound is monotone by inclusion
(\cref{prop:restrict-principle}), so minimizing the gap selects the largest
family; an arbitrary replacement joint supplies no common evidence
(\cref{prop:restrict-unrelated-joint-elbo}); and evidence-preserving channels,
which do order bounds across latent spaces, do so without selecting a
nondegenerate partition (\cref{sec:cg-elbo}). The problem is settled
positively by a functional on admissible channels or partitions, expressible
relative to one fixed joint and not monotone toward an endpoint, whose
optimizer is nondegenerate, and negatively by a proof of endpoint degeneracy
for a precisely stated class of such functionals.

\section{The factorwise decomposition}
\label{sec:elbo-factorization}

\propositionheading{Exact factorwise decomposition}{prop:elbo-factorwise-decomposition}
\label{hyp:elbo-factorwise-integrability}
Assume (H1)--(H3) of \cref{hyp:elbo-evidence-domain} and factorwise
log-integrability: the absolute logarithm of every component density of
\eqref{eq:gen-density} and every record energy $E_{a,o}$ is
$\mathbb Q_{V,o,X}$-integrable. Then (i) the generative half
\eqref{eq:elbo-log-integrability} of (H4) holds, and (H4) holds exactly when the
recognition half \eqref{eq:prob-integrability} does; (ii) under (H4) the free
energy has exactly one decomposition into factor terms,
\begin{align}
\Fenergy_V[\mathbb Q_{V,o,X};X,o]
={}&\sum_{i\in V}\E_{\mathbb Q_{V,o,X}}\bigl[-\log p^m_{\theta,i}(m_i\given m_{\operatorname{pa}(i)},X)\bigr]
+\sum_{i\in V}\E_{\mathbb Q_{V,o,X}}\bigl[-\log p^k_{\theta,i}(k_i\given k_{\operatorname{pa}(i)},m_i,X)\bigr]
\notag\\
&+\sum_{a\in\mathfrak A}\E_{\mathbb Q_{V,o,X}}\bigl[E_{a,o}(Y_{\partial a})\bigr]
+\E_{\mathbb Q_{V,o,X}}\bigl[\log q_{V,o,X}(Y)\bigr],
\label{eq:elbo-factorization}
\end{align}
with $\operatorname{pa}(i)=\varnothing$ at the roots, $2N+|\mathfrak A|$
generative factors, and a final term that is one full-joint recognition log
density, not a sum of marginal entropies; (iii) factorwise integrability does
not imply the recognition half, and the generative half does not imply
factorwise integrability. \status{ESTABLISHED}

\noindent\emph{Proof.}
By \eqref{eq:gen-density} and \cref{hyp:local-interaction-kernels},
\begin{equation}
\log p_{\theta,V}(o,y\given X)
=\sum_{i\in V}\bigl[\log p^m_{\theta,i}(m_i\given m_{\operatorname{pa}(i)},X)
+\log p^k_{\theta,i}(k_i\given k_{\operatorname{pa}(i)},m_i,X)\bigr]
-\sum_{a\in\mathfrak A}E_{a,o}(y_{\partial a}),
\label{eq:elbo-audit}
\end{equation}
a finite sum, so the triangle inequality gives (i). Under (H4), taking
$\mathbb Q_{V,o,X}$-expectations of the negative of \eqref{eq:elbo-audit}
term by term and adding $\E_{\mathbb Q_{V,o,X}}[\log q_{V,o,X}]$ gives
\eqref{eq:elbo-factorization} from \eqref{eq:elbo-definition}. For (iii),
\cref{sec:app-integrability-witnesses} gives a model with bounded factor
logarithms whose recognition density has an infinite mean logarithm, and a
root law $\mathrm c/(m(\log m)^2)$ whose state bridge cancels its logarithm, so
that the complete log joint is integrable while one factor is not.
\hfill$\square$

\section{Observations as interaction records}
\label{sec:local-observation-equivalence}

\definitionheading{Operational environment node}{def:obs-operational-environment-node}
An operational environment node is a standard Borel state space with a
normalized state law and a measurable message kernel.

\theoremheading{Observation--interaction kernel equivalence}{thm:obs-agent-interaction-equivalence}
Every normalized observation kernel $K(do\given y)$ from a measurable space to a
standard Borel space is the message of an operational environment node: there
are $U\sim\operatorname{Unif}[0,1]$ and a measurable map $F$ such that
\begin{equation}
O=F(Y,U),
\qquad
\Pr\bigl(F(y,U)\in A\bigr)=K(A\given y).
\label{eq:obs-randomization}
\end{equation}
Conversely, marginalizing any environment state and message policy produces a
normalized observation kernel. \status{ESTABLISHED}

\noindent\emph{Proof.}
The forward direction is the randomization lemma for kernels
\citep{Kallenberg2021}. By Kuratowski's theorem a standard Borel space admits a
Borel isomorphism $\phi$ onto a Borel subset of $\R$. The function
$G(y,x)=K(\phi^{-1}(-\infty,x]\given y)$ is measurable in $y$ and a
distribution function in $x$, so the quantile
$F_0(y,u)=\inf\{x:G(y,x)\geq u\}$ on $(0,1)$ satisfies
$\{F_0\leq x\}=\{(y,u):u\leq G(y,x)\}$; hence $F_0$ is jointly measurable and
$\Pr(F_0(y,U)\leq x)=G(y,x)$. The value $F_0(y,U)$ lies in the image of $\phi$
almost surely; setting $F=\phi^{-1}\circ F_0$ there and a fixed point elsewhere
gives \eqref{eq:obs-randomization}. Conversely, a state law
$\varrho_{\mathrm{env}}(du)$ and a message policy
$\kappa_{\mathrm{env}}(y,u,A)$ give
$K(A\given y)=\int\kappa_{\mathrm{env}}(y,u,A)\varrho_{\mathrm{env}}(du)$, which
is measurable in $y$, countably additive in $A$, and normalized.
\hfill$\square$

\propositionheading{Conditioning information and retained statistics}{prop:obs-conditioning-information}
Let $(Y,O)$ have joint law $P$ on standard Borel spaces with marginals $P^Y$ and
$P^O$, let $I(Y;O)=\KL(P\Vert P^Y\otimes P^O)$, and for a measurable statistic
$S=T(O)$ let $I(Y;O\given S)=\int\KL(P_{Y,O\given s}\Vert P_{Y\given s}\otimes P_{O\given s})P^S(ds)$.
The information carried by the record is
\begin{equation}
I(Y;O)=\int\KL\bigl(P(dY\given o)\Vert P^Y\bigr)P^O(do)\in[0,+\infty].
\label{eq:obs-conditioning-information}
\end{equation}
Moreover
\begin{equation}
I(Y;O)=I(Y;S)+I(Y;O\given S)
\quad\text{in }[0,+\infty],
\label{eq:obs-conditioning-statistic-chain}
\end{equation}
so replacing $O$ by $S$ leaves the posterior unchanged exactly when
$I(Y;O\given S)=0$, up to the choice of conditional versions. Deleting the record
makes $S$ constant and returns the prior; for
$Y\sim\operatorname{Ber}(1/2)$ and $O=Y$, $I(Y;O)=\log2$ and
$P(dY\given O=o)=\delta_o$. A kernel that does not depend on $y$ carries zero
information. \status{ESTABLISHED}

\noindent\emph{Proof.}
Disintegrating $\KL(P\Vert P^Y\otimes P^O)$ over $O$ with
\cref{lem:prob-kl-chain-rule}, whose $O$-marginals agree, gives
\eqref{eq:obs-conditioning-information}. Since $o\mapsto(T(o),o)$ is a
bimeasurable bijection onto its graph, relative entropy is unchanged by passing
to $(Y,S,O)$ \eqref{eq:cg-kl-bimeasurable-invariance}. Disintegrating
$\KL(P_{Y,S,O}\Vert P^Y\otimes P_{S,O})$ over $(Y,S)$ gives
$I(Y;S)+\E_{P_{Y,S}}\KL(P_{O\given Y,S}\Vert P_{O\given S})$, because under the
reference $O$ is independent of $Y$ given $S$ with law $P_{O\given S}$; for
each $s$, disintegrating
$\KL(P_{Y,O\given s}\Vert P_{Y\given s}\otimes P_{O\given s})$ over $Y$ and
integrating in $s$ identifies the second term with $I(Y;O\given S)$, which
proves \eqref{eq:obs-conditioning-statistic-chain}; the discrete case is
\citet[Sec.~2.5]{CoverThomas2006}. Posterior sufficiency is the vanishing of
the conditional term, the Bernoulli values are direct, and a $y$-independent
kernel makes $P=P^Y\otimes P^O$. \hfill$\square$

\remarkheading{Environment nodes in the variational problem}{rem:obs-environment-fiber}
\Cref{thm:obs-agent-interaction-equivalence} is exact for the joint law of
$(Y,O)$, and it keeps the record variable $O$ and hence the conditioning
$\sigma$-algebra $\sigma(O)$, but it is not exact for a variational problem
over the enlarged state $(Y,U)$. For
$P^O$-almost every $o$ the posterior of $(Y,U)$ is concentrated on the level set
$\{F(y,u)=o\}$. For $Y\sim\mathcal N(0,1)$ and
$F(y,u)=y+F_{\mathcal N}^{-1}(u)$, with $F_{\mathcal N}$ the standard normal
distribution function, which realizes $K(do\given y)=\mathcal N(y,1)$, that
level set is the Lebesgue-null curve $\{u=F_{\mathcal N}(o-y)\}$, so every
recognition law on $(Y,U)$ with a Lebesgue density has relative entropy
$+\infty$ by \cref{prop:elbo-subspace-support-singular}, whereas the
fiber-restricted law $q(dy)\delta_{F_{\mathcal N}(o-y)}(du)$, the image of $q$
under the injective map $y\mapsto(y,F_{\mathcal N}(o-y))$ that also carries the
posterior of $Y$ to that of $(Y,U)$, has relative entropy
$\KL(q\Vert P(dY\given o))$. A finite variational problem therefore integrates
$U$ out or restricts recognition to that fiber. \status{ESTABLISHED}

\section{Local conditional free energy and the potential identity}
\label{sec:local-conditional-vfe}

Fix a nonempty block $B\subseteq V$, one agent or any aggregate, and write
$y=(y_B,y_{-B})$; the probability model does not change with $B$.

\lemmaheading{Block disintegration of the posterior}{lem:obs-block-posterior}
Let $\mathfrak A_B=\{a\in\mathfrak A:\partial a\cap B\ne\varnothing\}$ and define
the incident and outside likelihoods and energies
\begin{align}
L_{B,o}(y_B;y_{-B})&=\prod_{a\in\mathfrak A_B}e^{-E_{a,o}(y_{\partial a})},
&L_{\bar B,o}(y_{-B})&=\prod_{a\notin\mathfrak A_B}e^{-E_{a,o}(y_{\partial a})},
\label{eq:obs-block-incident-likelihood}\\
H_{B,o}(y_B;y_{-B})&=-\log L_{B,o}=\sum_{a\in\mathfrak A_B}E_{a,o}(y_{\partial a}),
&H_{\bar B,o}(y_{-B})&=-\log L_{\bar B,o}(y_{-B}),
\label{eq:obs-block-incident-energy}
\end{align}
so that $L_o=L_{B,o}L_{\bar B,o}$ and $H_o=H_{B,o}+H_{\bar B,o}$ in
$(-\infty,+\infty]$. Fix a measurable version of the baseline disintegration
\begin{equation}
\mathbb P^{\mathsf Y}_{\theta,V}(dy_B,dy_{-B})
=\mathbb P^{\mathsf Y}_{\theta,V,-B}(dy_{-B})\mathbb P^{\mathsf Y}_{\theta,V,B}(dy_B\given y_{-B}),
\label{eq:obs-baseline-block-disintegration}
\end{equation}
and define the local evidence and the likelihood of the outside state,
\begin{align}
Z_B(y_{-B})&=\int L_{B,o}(y_B;y_{-B})\mathbb P^{\mathsf Y}_{\theta,V,B}(dy_B\given y_{-B}),
\label{eq:obs-local-evidence}\\
\widetilde L_{B,o}(y_{-B})&=\int L_o(y_B,y_{-B})\mathbb P^{\mathsf Y}_{\theta,V,B}(dy_B\given y_{-B}).
\label{eq:obs-full-complement-likelihood}
\end{align}
Then (i) the evidence and the outside posterior marginal are
\begin{equation}
p_{\theta,V}(o\given X)=\int\widetilde L_{B,o}d\mathbb P^{\mathsf Y}_{\theta,V,-B},
\qquad
\boldsymbol\Pi_{\theta,V,o,X,-B}(dy_{-B})
=p_{\theta,V}(o\given X)^{-1}\widetilde L_{B,o}(y_{-B})\mathbb P^{\mathsf Y}_{\theta,V,-B}(dy_{-B});
\label{eq:obs-posterior-outside}
\end{equation}
(ii) the set $\Omega_{B,o}$ of outside states at which
$0<\widetilde L_{B,o}<\infty$, $0<L_{\bar B,o}<\infty$, $0<Z_B<\infty$, and
$\widetilde L_{B,o}=L_{\bar B,o}Z_B$ has $\boldsymbol\Pi_{\theta,V,o,X,-B}$-measure
one, though it need not be $\mathbb P^{\mathsf Y}_{\theta,V,-B}$-full; (iii) the
kernel
\begin{equation}
\boldsymbol\Pi_{\theta,V,o,X,B}(dy_B\given y_{-B})
=Z_B(y_{-B})^{-1}L_{B,o}(y_B;y_{-B})\mathbb P^{\mathsf Y}_{\theta,V,B}(dy_B\given y_{-B})
\qquad(y_{-B}\in\Omega_{B,o}),
\label{eq:obs-local-posterior}
\end{equation}
extended by any measurable probability kernel off $\Omega_{B,o}$, satisfies
\begin{equation}
\boldsymbol\Pi_{\theta,V,o,X}(dy_B,dy_{-B})
=\boldsymbol\Pi_{\theta,V,o,X,-B}(dy_{-B})\boldsymbol\Pi_{\theta,V,o,X,B}(dy_B\given y_{-B}).
\label{eq:obs-posterior-block-disintegration}
\end{equation}
The posterior determines this kernel, and every pointwise local quantity built
from it, only $\boldsymbol\Pi_{\theta,V,o,X,-B}$-almost everywhere; two baseline
versions give local posteriors that agree there. \status{ESTABLISHED}

\noindent\emph{Proof.}
Tonelli's theorem applied to
\eqref{eq:obs-global-evidence}--\eqref{eq:obs-global-posterior} with
\eqref{eq:obs-baseline-block-disintegration} gives (i); defining $\widetilde L_{B,o}$ directly avoids the product
$L_{\bar B,o}Z_B$, which can be $0\cdot\infty$ off $\Omega_{B,o}$. For (ii),
finite evidence makes $\widetilde L_{B,o}<\infty$ off a
$\mathbb P^{\mathsf Y}_{\theta,V,-B}$-null set, and by (i) the outside posterior
charges only $\{\widetilde L_{B,o}>0\}$. There $L_{\bar B,o}$ is finite because
each factor density is, positive because $L_{\bar B,o}=0$ would force
$\widetilde L_{B,o}=0$, and independent of $y_B$, so it leaves the conditional
integral and $Z_B=\widetilde L_{B,o}/L_{\bar B,o}\in(0,\infty)$. For (iii), a
bounded measurable $f$ has
$\int fd\boldsymbol\Pi_{\theta,V,o,X}
=p_{\theta,V}(o\given X)^{-1}\int L_{\bar B,o}\int fL_{B,o}d\mathbb P^{\mathsf Y}_{\theta,V,B}d\mathbb P^{\mathsf Y}_{\theta,V,-B}$,
and on $\Omega_{B,o}$ the inner integral is $Z_B$ times the integral of $f$
against \eqref{eq:obs-local-posterior}; (i) and (ii) turn the outer weight into
$\boldsymbol\Pi_{\theta,V,o,X,-B}$. Almost-everywhere uniqueness is uniqueness
of disintegrations; two baseline versions agree
$\mathbb P^{\mathsf Y}_{\theta,V,-B}$-almost everywhere, hence
$\boldsymbol\Pi_{\theta,V,o,X,-B}$-almost everywhere by (i). \hfill$\square$

\definitionheading{Local conditional free energy}{thm:obs-local-multiagent-elbo}
For $y_{-B}\in\Omega_{B,o}$ and every probability law $r_B(\cdot\given y_{-B})$
on $\mathsf Y_B$, the local conditional ELBO and free energy of the block are
\begin{equation}
\Fenergy^{\mathrm{ext}}_{B,o}(r_B;y_{-B}):=-\Lelbo^{\mathrm{ext}}_{B,o}(r_B;y_{-B})
:=-\log Z_B(y_{-B})
+\KL\bigl(r_B(\cdot\given y_{-B})\Vert\boldsymbol\Pi_{\theta,V,o,X,B}(\cdot\given y_{-B})\bigr).
\label{eq:obs-local-vfe}
\end{equation}
For $B=\{i\}$ this is the local free energy of agent $i$. It contains every
record incident to the block and no other.

\corollaryheading{Local Gibbs bound and the exact block E-coordinate}{prop:elbo-exact-e-coordinate}
(i) $\Lelbo^{\mathrm{ext}}_{B,o}(r_B;y_{-B})\leq\log Z_B(y_{-B})$, with
equality exactly at $r_B=\boldsymbol\Pi_{\theta,V,o,X,B}(\cdot\given y_{-B})$.
If $r_B\ll\mathbb P^{\mathsf Y}_{\theta,V,B}(\cdot\given y_{-B})$ with finite
relative entropy and the negative part of $H_{B,o}(\cdot;y_{-B})$ is
$r_B$-integrable,
\begin{equation}
\Fenergy^{\mathrm{ext}}_{B,o}(r_B;y_{-B})
=\KL\bigl(r_B(\cdot\given y_{-B})\Vert\mathbb P^{\mathsf Y}_{\theta,V,B}(\cdot\given y_{-B})\bigr)
+\E_{r_B(\cdot\given y_{-B})}H_{B,o}(Y_B;y_{-B}).
\label{eq:obs-local-energy-form}
\end{equation}

(ii) If $Q(dy_B,dy_{-B})=Q_{-B}(dy_{-B})r_B(dy_B\given y_{-B})$ and
$Q_{-B}\ll\boldsymbol\Pi_{\theta,V,o,X,-B}$, then in $[0,+\infty]$, with the
baseline conditional of \eqref{eq:obs-baseline-block-disintegration} in the
second identity,
\begin{align}
\KL\bigl(Q\Vert\boldsymbol\Pi_{\theta,V,o,X}\bigr)
&=\KL\bigl(Q_{-B}\Vert\boldsymbol\Pi_{\theta,V,o,X,-B}\bigr)
+\E_{Q_{-B}}\KL\bigl(r_B(\cdot\given Y_{-B})\Vert\boldsymbol\Pi_{\theta,V,o,X,B}(\cdot\given Y_{-B})\bigr),
\label{eq:obs-posterior-block-kl-chain}\\
\KL\bigl(Q\Vert\mathbb P^{\mathsf Y}_{\theta,V}\bigr)
&=\KL\bigl(Q_{-B}\Vert\mathbb P^{\mathsf Y}_{\theta,V,-B}\bigr)
+\E_{Q_{-B}}\KL\bigl(r_B(\cdot\given Y_{-B})\Vert\mathbb P^{\mathsf Y}_{\theta,V,B}(\cdot\given Y_{-B})\bigr).
\label{eq:obs-block-kl-chain}
\end{align}

(iii) Let $\mathcal B$ be any block of latent coordinates with complement
$-\mathcal B$, let $P_-$ be the $-\mathcal B$ marginal of
$P=\boldsymbol\Pi_{\theta,V,o,X}$, let $R$ be any law on the $-\mathcal B$
coordinates, and let $\kappa_{\mathcal B}(dY_{\mathcal B}\given Y_{-\mathcal B})$
be a Markov kernel that agrees $P_-$-almost everywhere with a regular
conditional law of $P$ (for an agent block, \eqref{eq:obs-local-posterior}).
Among all laws with $-\mathcal B$ marginal $R$, the divergence from $P$ has
minimum $\KL(R\Vert P_-)\in[0,+\infty]$, attained by
\begin{equation}
\mathbb Q^{\mathrm{new}}_{V,o,X}(dY_{\mathcal B},dY_{-\mathcal B})
=\kappa_{\mathcal B}(dY_{\mathcal B}\given Y_{-\mathcal B})R(dY_{-\mathcal B}),
\label{eq:elbo-e-coordinate}
\end{equation}
which therefore maximizes $\Lelbo_V^{\mathrm{ext}}$ on that class. The minimizer
is unique when $\KL(R\Vert P_-)<\infty$; otherwise every law in the class has
extended bound $-\infty$. If the minimum is finite and
\begin{equation}
\E_{\mathbb Q^{\mathrm{new}}_{V,o,X}}\bigl[\lvert\log p_{\theta,V}(o,Y\given X)\rvert
+\lvert\log q^{\mathrm{new}}_{V,o,X}(Y)\rvert\bigr]<\infty,
\label{eq:elbo-e-coordinate-h4}
\end{equation}
the new law satisfies (H3)--(H4) and maximizes the classical $\Lelbo_V$ among
the laws with marginal $R$ that satisfy \cref{hyp:elbo-evidence-domain}.
Finite complement relative entropy does not imply \eqref{eq:elbo-e-coordinate-h4}.

(iv) If $\Fenergy^{\mathrm{ext}}_o(Q)<\infty$ for $Q=Q_{-B}r_B$, then for
$\mathrm w\in[0,1]$ the damped kernel
$r_B^{\mathrm w}=(1-\mathrm w)r_B+\mathrm w\boldsymbol\Pi_{\theta,V,o,X,B}$
satisfies
\begin{equation}
\Fenergy^{\mathrm{ext}}_o(Q_{-B}r_B^{\mathrm w})
\leq(1-\mathrm w)\Fenergy^{\mathrm{ext}}_o(Q)
+\mathrm w\Fenergy^{\mathrm{ext}}_o(Q_{-B}\boldsymbol\Pi_{\theta,V,o,X,B})
\leq\Fenergy^{\mathrm{ext}}_o(Q),
\label{eq:obs-damped-block-update}
\end{equation}
so replacing $r_B$ by the local posterior with $Q_{-B}$ fixed does not increase
$\Fenergy^{\mathrm{ext}}_o$. \status{ESTABLISHED}

\noindent\emph{Proof.}
For (i), the conditional model with reference
$\mathbb P^{\mathsf Y}_{\theta,V,B}(\cdot\given y_{-B})$ and likelihood
$L_{B,o}(\cdot;y_{-B})$ has evidence $Z_B(y_{-B})\in(0,\infty)$ and posterior
\eqref{eq:obs-local-posterior}, and \cref{thm:elbo-extended-gap} applies to it.
Part (ii) is \cref{lem:prob-kl-chain-rule} with the disintegrations
\eqref{eq:obs-posterior-block-disintegration} and
\eqref{eq:obs-baseline-block-disintegration}. For (iii), the fixed-marginal
clause of \cref{lem:prob-kl-chain-rule}, with the $-\mathcal B$ coordinates
first, gives the minimum $\KL(R\Vert P_-)$, the minimizer
\eqref{eq:elbo-e-coordinate}, and uniqueness when the minimum is finite; then
$R\ll P_-$, the completion of $\kappa_{\mathcal B}$ off a $P_-$-null set is
irrelevant, and $d\mathbb Q^{\mathrm{new}}_{V,o,X}/dP=(dR/dP_-)(Y_{-\mathcal B})$.
When the minimum is infinite, the data-processing clause
\eqref{eq:cg-kl-dpi-extended} applied to the projection makes every divergence
in the class infinite. \Cref{def:elbo-extended} converts divergence into the
extended bound, and \cref{thm:elbo-exact-identity} gives the classical
statement. The countable witness of \cref{sec:app-integrability-witnesses} has
$R=P_-$ and fails \eqref{eq:elbo-e-coordinate-h4}. For (iv), finiteness makes
$Q_{-B}\ll\boldsymbol\Pi_{\theta,V,o,X,-B}$ by data processing, and
\eqref{eq:obs-posterior-block-kl-chain} leaves the outside term unchanged; the
exact replacement sets the conditional term to zero, and convexity of relative
entropy in its first argument, applied for each $y_{-B}$, gives
\eqref{eq:obs-damped-block-update}. \hfill$\square$

The exact full E-step is case (iii) with $\mathcal B$ equal to all
coordinates; a coordinate that optimizes inside a product family does not have
the form \eqref{eq:elbo-e-coordinate} and leaves the gap of
\cref{prop:restrict-principle}. By (iv) a sequential sweep of exact or damped
block updates does not increase the collective free energy; replacing the
conditionals of several correlated blocks at once is not covered, because each
replacement changes the outside marginal of the others.

\theoremheading{Exact local--global potential identity}{thm:obs-local-global-potential}
Let $Q=Q_{-B}r_B$ and $Q'=Q_{-B}r'_B$ have the same outside marginal with
$Q_{-B}\ll\boldsymbol\Pi_{\theta,V,o,X,-B}$, and suppose both have finite
divergence from $\boldsymbol\Pi_{\theta,V,o,X}$. Then
\begin{equation}
\Fenergy^{\mathrm{ext}}_o(Q')-\Fenergy^{\mathrm{ext}}_o(Q)
=\E_{Q_{-B}}\bigl[\Fenergy^{\mathrm{ext}}_{B,o}(r'_B;Y_{-B})
-\Fenergy^{\mathrm{ext}}_{B,o}(r_B;Y_{-B})\bigr].
\label{eq:obs-local-global-potential}
\end{equation}
When in addition every baseline relative entropy and the positive and negative
parts of the energies are integrable,
\begin{equation}
\Fenergy^{\mathrm{ext}}_o(Q)
=\KL\bigl(Q_{-B}\Vert\mathbb P^{\mathsf Y}_{\theta,V,-B}\bigr)
+\E_{Q_{-B}}\bigl[H_{\bar B,o}(Y_{-B})+\Fenergy^{\mathrm{ext}}_{B,o}(r_B;Y_{-B})\bigr].
\label{eq:obs-local-global-decomposition}
\end{equation}
Every exact local update is thus a coordinate update of the collective free
energy. The fixed outside marginal is load bearing. \status{ESTABLISHED}

\noindent\emph{Proof.}
Apply \eqref{eq:obs-posterior-block-kl-chain} to both laws. The finite
divergences make both conditional terms integrable, the outside terms agree,
the constant $-\log p_{\theta,V}(o\given X)$ cancels, and so does
$-\log Z_B(Y_{-B})$ inside the two local functionals. The decomposition follows
from \eqref{eq:obs-block-kl-chain}, \eqref{eq:obs-local-energy-form}, and
\eqref{eq:obs-collective-energy-form}. \hfill$\square$

\corollaryheading{Exact potential game}{cor:obs-potential-game}
Let the blocks be singletons and choose a nonempty strategy set $\mathcal R_i$
for each agent. Suppose every profile $(r_i)\in\prod_{i\in V}\mathcal R_i$
defines a product recognition law $Q=\bigotimes_{i\in V}r_i$ with
$\KL(Q\Vert\boldsymbol\Pi_{\theta,V,o,X})<\infty$ and
$\E_{Q_{-i}}\lvert\log Z_{\{i\}}(Y_{-i})\rvert<\infty$ for every $i$, where
$Q_{-i}=\bigotimes_{j\ne i}r_j$. The outside-averaged local free energy
\begin{equation}
\overline\Fenergy_i(r_i;r_{-i}):=\E_{Q_{-i}}\Fenergy^{\mathrm{ext}}_{\{i\},o}(r_i;Y_{-i})
\label{eq:obs-outside-averaged-local-vfe}
\end{equation}
is finite, and between two such profiles that differ only in $r_i$ it changes by
exactly the change of $\Fenergy^{\mathrm{ext}}_o$. The players $V$, strategy sets
$\mathcal R_i$, and costs $\overline\Fenergy_i$ therefore form an exact potential game
with potential $\Fenergy^{\mathrm{ext}}_o$ in the sense of
\citet[Sec.~2]{MondererShapley1996}, written for costs instead of payoffs.
\status{ESTABLISHED}

\noindent\emph{Proof.}
Finite divergence gives $Q_{-i}\ll\boldsymbol\Pi_{\theta,V,o,X,-i}$ by data
processing, and \eqref{eq:obs-posterior-block-kl-chain} makes the averaged
conditional divergence finite; with the hypothesis on $\log Z_{\{i\}}$ this
makes \eqref{eq:obs-outside-averaged-local-vfe} finite. A unilateral change
leaves $Q_{-i}$ fixed, and \cref{thm:obs-local-global-potential} with
$B=\{i\}$ gives the equality of changes. \hfill$\square$

\Cref{fig:obs-local-global-potential} shows which records enter a block
potential; block-gradient dynamics are in \cref{sec:local-natural-gradient}.

\begin{figure}[htbp]
\centering
\begin{tikzpicture}[>=stealth,scale=0.85,every node/.style={transform shape}]
  \node[draw,circle,fill=lightblue,minimum size=0.85cm] (a1) at (0,1.5) {$1$};
  \node[draw,circle,fill=lightblue,minimum size=0.85cm] (a2) at (2.1,2.3) {$2$};
  \node[draw,circle,fill=lightgreen,minimum size=0.85cm] (a3) at (4.2,1.5) {$3$};
  \node[draw,circle,fill=lightgreen,minimum size=0.85cm] (a4) at (5.2,-0.2) {$4$};
  \node[draw,circle,fill=lightblue,minimum size=0.85cm] (a5) at (1.2,-0.5) {$5$};
  \node[draw,rectangle,fill=lightorange,minimum size=0.52cm] (f12) at (1.05,1.9) {$o_{a}$};
  \node[draw,rectangle,fill=lightorange,minimum size=0.52cm] (f235) at (2.65,0.45) {$o_{a'}$};
  \node[draw,rectangle,fill=lightorange,minimum size=0.52cm] (f34) at (4.7,0.65) {$o_{a''}$};
  \draw (a1)--(f12)--(a2);
  \draw (a2)--(f235)--(a3);
  \draw (a5)--(f235);
  \draw (a3)--(f34)--(a4);
  \draw[rounded corners,very thick,sciencegreen]
    (3.55,2.15) rectangle (5.75,-0.85);
  \node[text=sciencegreen] at (4.65,2.48) {block $B=\{3,4\}$};
  \draw[->,very thick,primaryblue] (6.05,0.65)--(7.35,0.65);
  \node[align=center,draw,rounded corners,fill=lightblue,text width=3.5cm]
    at (9.25,0.65)
    {$\Delta\Fenergy_o^{\rm ext}
      =\E_{Q_{-B}}\Delta\Fenergy_{B,o}^{\rm ext}$\\
     incident records: $o_{a'},o_{a''}$};
\end{tikzpicture}
\caption{A block-local free energy contains every record incident to the block:
$o_{a'}$ crosses the boundary and enters the block conditional, while $o_a$ lies
outside and cancels.}
\label{fig:obs-local-global-potential}
\end{figure}

\section{Additive accounting}
\label{sec:local-additive-accounting}

The local free energies of \cref{thm:obs-local-multiagent-elbo} are coordinate
potentials of the collective free energy, not shares of it; an additive
single-count ledger is a different construction.

\corollaryheading{Correlated ledger and incident overcounting}{cor:obs-correlated-ledger}
Suppose $\mathbb P^{\mathsf Y}_{\theta,V}=\bigotimes_{i\in V}\varrho_i$ and let
$Q$ be any law on the latent space. Then (i) in $[0,+\infty]$,
\begin{equation}
\KL\bigl(Q\Vert\mathbb P^{\mathsf Y}_{\theta,V}\bigr)
=\operatorname{TC}(Q)+\sum_i\KL(Q_i\Vert\varrho_i);
\label{eq:obs-correlated-complexity-ledger}
\end{equation}
(ii) whenever \eqref{eq:obs-collective-energy-form} holds,
\begin{equation}
\Fenergy^{\mathrm{ext}}_o(Q)
=\operatorname{TC}(Q)+\sum_i\KL(Q_i\Vert\varrho_i)
+\E_Q\Bigl[\sum_{a\in\mathfrak A}E_{a,o}(Y_{\partial a})\Bigr],
\label{eq:obs-global-ledger}
\end{equation}
in which each record appears once and $\operatorname{TC}(Q)=0$ exactly when $Q$
is a product law, so mean field is a special case of the ledger and not a
premise of it; (iii) if every record has a nonempty scope, the singleton
incident energies satisfy, pointwise in $(-\infty,+\infty]$,
\begin{equation}
H_{\{i\},o}(y_i;y_{-i})=\sum_{a:i\in\partial a}E_{a,o}(y_{\partial a}),
\qquad
\sum_{i\in V}H_{\{i\},o}(y_i;y_{-i})=\sum_{a\in\mathfrak A}\lvert\partial a\rvert E_{a,o}(y_{\partial a}),
\label{eq:obs-singleton-incident-counting}
\end{equation}
so the singleton local free energies count each record $\lvert\partial a\rvert$
times and are not the single-count ledger; (iv) this is a statement about these
local free energies, not a bar to additive decompositions: iterating
\cref{lem:prob-kl-chain-rule} along any ordering of $V$ writes
$\KL(Q\Vert\boldsymbol\Pi_{\theta,V,o,X})$ as one term per agent, each
conditioned on the agents before it. \status{ESTABLISHED}

\noindent\emph{Proof.}
Part (i) is \cref{thm:elbo-total-correlation-chain} with
$\varrho_j=\varrho_i$, and (ii) substitutes it into
\eqref{eq:obs-collective-energy-form}. For (iii), every $E_{a,o}$ lies in
$(-\infty,+\infty]$ and both index sets are finite, so the rearrangement is
valid pointwise. Part (iv) is repeated application of the chain rule, peeling
one agent at a time. \hfill$\square$

The same chain rule accounts for link values once they are latent; the
interaction multigraph itself stays structural (\cref{sec:intro-conventions}).

\propositionheading{Collective free energy with random link values}{prop:elbo-random-link-vfe}
Let the link family $\Theta=(\Theta^x_e)$ on the edge copies of $\mathfrak I$ be
a latent variable with values in a standard Borel space $\mathscr T$. Let
$\nu^{\mathscr T}$ be a probability law on $\mathscr T$, let
$\mathscr S:\mathscr T\to(-\infty,+\infty]$ be a measurable action with
$0<Z_{\mathscr S}=\int e^{-\mathscr S}d\nu^{\mathscr T}<\infty$, and let the link
prior be $Z_{\mathscr S}^{-1}e^{-\mathscr S}\nu^{\mathscr T}$. For each $\Theta$
let the record model of \cref{hyp:local-interaction-kernels} with links $\Theta$
have evidence $p(o\given\Theta)\in(0,\infty)$, measurable in $\Theta$, posterior
$\boldsymbol\Pi_\Theta$, and free energy
$\Fenergy_{\mathrm{matter}}(Q_\Theta;o,\Theta)=-\log p(o\given\Theta)+\KL(Q_\Theta\Vert\boldsymbol\Pi_\Theta)$,
and assume the joint evidence
$p(o)=\int p(o\given\Theta)Z_{\mathscr S}^{-1}e^{-\mathscr S(\Theta)}\nu^{\mathscr T}(d\Theta)$
is finite. For a recognition law $Q(d\Theta,dy)=R(d\Theta)Q_\Theta(dy)$ with
$\KL(R\Vert\nu^{\mathscr T})<\infty$ and $\E_R\lvert\mathscr S\rvert<\infty$, the
extended free energy of the joint model is
\begin{equation}
\Fenergy=\E_R\Fenergy_{\mathrm{matter}}(Q_\Theta;o,\Theta)+\E_R\mathscr S
+\KL(R\Vert\nu^{\mathscr T})+\log Z_{\mathscr S}
\quad\text{in }(-\infty,+\infty].
\label{eq:rg-geometry-random-link-vfe}
\end{equation}
Omitting the link entropy term or the normalizer changes the free energy; when
the links are structural parameters instead of latent variables, only
$\Fenergy_{\mathrm{matter}}$ appears. \status{ESTABLISHED}

\noindent\emph{Proof.}
The posterior of the joint model has conditional law $\boldsymbol\Pi_\Theta$
given $\Theta$ and link marginal
$\boldsymbol\Pi^{\mathscr T}(d\Theta)
=p(o)^{-1}p(o\given\Theta)Z_{\mathscr S}^{-1}e^{-\mathscr S(\Theta)}\nu^{\mathscr T}(d\Theta)$,
so the chain rule
\eqref{eq:elbo-kl-chain-rule} gives
$\KL(Q\Vert\boldsymbol\Pi)=\KL(R\Vert\boldsymbol\Pi^{\mathscr T})+\E_R\KL(Q_\Theta\Vert\boldsymbol\Pi_\Theta)$
in $[0,+\infty]$, with a measurable integrand. Since
$\E_R\lvert\mathscr S\rvert<\infty$, $R$ charges no set where
$\mathscr S=+\infty$, so $R\ll\boldsymbol\Pi^{\mathscr T}$ and
$\log(dR/d\boldsymbol\Pi^{\mathscr T})=\log(dR/d\nu^{\mathscr T})+\mathscr S+\log Z_{\mathscr S}
-\log p(o\given\Theta)+\log p(o)$. The first two terms are $R$-integrable, and so
is the positive part $f$ of $\log p(o\given\cdot)$: for $f_n=\min(f,n)$, the
nonnegative divergence from the law proportional to $e^{f_n}$ times the prior
gives $\E_Rf_n\leq\KL(R\Vert Z_{\mathscr S}^{-1}e^{-\mathscr S}\nu^{\mathscr T})+\log(1+p(o))$,
and that divergence equals
$\KL(R\Vert\nu^{\mathscr T})+\E_R\mathscr S+\log Z_{\mathscr S}$. Integrating
term by term and adding the conditional term, whose integrand is nonnegative,
turns $\Fenergy=-\log p(o)+\KL(Q\Vert\boldsymbol\Pi)$ into
\eqref{eq:rg-geometry-random-link-vfe}; every summand has an integrable negative
part, so no difference of infinities occurs. \hfill$\square$

\Cref{ch:links} supplies, for every Lie group $G$, data that satisfy the hypotheses of this proposition:
the latent variable is the family of Lie-algebra generators of the links, the base law is Gaussian and
the action a nonnegative combination of loop defects, the joint has unit mass on every record hypergraph
(\Cref{prop:link-latent-normalization}), and for a Gaussian recognition law of the generators the link
sector of
\eqref{eq:rg-geometry-random-link-vfe} is the relative entropy from the generator prior
(\Cref{prop:link-gaussian-recognition-vfe}).

\section{Model coordinates and accepted updates}
\label{sec:elbo-coordinates}

In the two-argument form $\Fenergy_V(\mathbb Q;\theta)$ of
\cref{req:gen-typing-prohibition}, block coordinates move $\mathbb Q$ at fixed
$\theta$ and a model coordinate moves $\theta$ at fixed $\mathbb Q$; each
accepted value defines a new fixed joint.

\propositionheading{Exact M-coordinate}{prop:elbo-exact-m-coordinate}
Hold $\mathbb Q_{V,o,X}$ fixed, and suppose the parameterized family supplies a
common regular-observation set or jointly measurable pointwise versions for the
candidates compared, with
$\E_{\mathbb Q_{V,o,X}}\lvert\log q_{V,o,X}(Y)\rvert<\infty$ and
$\E_{\mathbb Q_{V,o,X}}\lvert\log p_{\theta',V}(o,Y\given X)\rvert<\infty$
for every candidate $\theta'$. If a maximizer exists in the declared parameter
space, then every
\begin{equation}
\theta^{\mathrm{new}}\in\operatorname*{argmax}_{\theta'}
\E_{\mathbb Q_{V,o,X}}\bigl[\log p_{\theta',V}(o,Y\given X)\bigr]
\label{eq:elbo-m-coordinate}
\end{equation}
maximizes $\Lelbo_{V,\theta'}(\mathbb Q_{V,o,X};X,o)$ over the same space. The
finiteness hypotheses cannot be dropped. \status{ESTABLISHED}

\noindent\emph{Proof.}
The recognition term $\E_{\mathbb Q_{V,o,X}}[\log q_{V,o,X}]$ is finite and does
not depend on $\theta'$. For the last sentence, on $\{2,3,\ldots\}$ let
$q_n=\mathrm c/(n(\log n)^2)$ with $\mathrm c$ normalizing,
$h_n=\mathbf 1_{\{n=2\}}-q_2$, and $r_n=1+\mathrm c_2h_n$ with
$0<\mathrm c_2<1/q_2$, so that $r$ is bounded, strictly positive, nonconstant,
and $\sum_nq_nr_n=1$. With a trivial record, let the two candidates have latent
laws $p_0=q$ and $p_1=qr$, and let the recognition law be $q$. The entropy of
$q$ is infinite (\cref{sec:app-integrability-witnesses}), so
$\E_q\log p_0=\E_q\log p_1=-\infty$ and both candidates maximize the expected
complete log-likelihood, whereas their extended bounds are
$\E_q\log(p_0/q)=0$ and $\E_q\log r=-\KL(q\Vert p_1)<0$. \hfill$\square$

\definitionheading{Generalized M-coordinate acceptance}{def:elbo-generalized-m-acceptance}
A generalized M-coordinate need not attain a maximum, but it must pass the
acceptance condition
\begin{equation}
\Lelbo^{\mathrm{ext}}_{V,\theta^{\mathrm{new}}}(\mathbb Q_{V,o,X};X,o)
\geq\Lelbo^{\mathrm{ext}}_{V,\theta^{\mathrm{old}}}(\mathbb Q_{V,o,X};X,o),
\label{eq:elbo-acceptance}
\end{equation}
evaluated at an observation in the regular sets of both candidates, or with
jointly measurable versions declared as part of the parameterized model. No
recognition density and no (H4) is required; on the classical domain the
condition is the same inequality for \eqref{eq:elbo-definition}, and
\eqref{eq:elbo-m-coordinate} is a sufficient condition.

\corollaryheading{Evidence-free acceptance test}{cor:elbo-acceptance-evidence-free}
Suppose $\mathbb Q_{V,o,X}\ll M_{\theta^{\mathrm{old}},V,o,X}$, the two slice
measures \eqref{eq:elbo-slice-measure} are mutually absolutely continuous, and
both extended values are finite. Then
\begin{equation}
\Lelbo^{\mathrm{ext}}_{V,\theta^{\mathrm{new}}}(\mathbb Q_{V,o,X};X,o)
-\Lelbo^{\mathrm{ext}}_{V,\theta^{\mathrm{old}}}(\mathbb Q_{V,o,X};X,o)
=\E_{\mathbb Q_{V,o,X}}\biggl[\log\frac{p_{\theta^{\mathrm{new}},V}(o,Y\given X)}{p_{\theta^{\mathrm{old}},V}(o,Y\given X)}\biggr],
\label{eq:elbo-acceptance-evidence-free}
\end{equation}
so the test \eqref{eq:elbo-acceptance} needs no evidence value.
\status{ESTABLISHED}

\noindent\emph{Proof.}
With $r_x=d\mathbb Q_{V,o,X}/dM_{\theta^x,V,o,X}$ for
$x\in\{\mathrm{old},\mathrm{new}\}$, the Radon--Nikodym chain rule gives
$\log r_{\mathrm{old}}-\log r_{\mathrm{new}}
=\log(dM_{\theta^{\mathrm{new}},V,o,X}/dM_{\theta^{\mathrm{old}},V,o,X})$ almost
surely, and this derivative is the ratio of the two joint densities at $o$. By
\cref{prop:elbo-relative-log-representation} the negative part of each
$\log r_x$ is integrable, so finiteness makes each $\log r_x$ integrable, and
subtracting the two representations \eqref{eq:elbo-relative-log} proves the
identity. \hfill$\square$

\propositionheading{Evidence monotonicity}{prop:elbo-evidence-monotonicity}
Suppose (H1)--(H2) hold for the fixed $\mathbb Q_{V,o,X}$ at both
$\theta^{\mathrm{old}}$ and $\theta^{\mathrm{new}}$. If the E-coordinate is exact
at the old parameter, $\mathbb Q_{V,o,X}=\boldsymbol\Pi_{\theta^{\mathrm{old}},V,o,X}$,
and the M-coordinate satisfies \eqref{eq:elbo-acceptance}, then
\begin{equation}
\log p_{\theta^{\mathrm{new}},V}(o\given X)
\geq\Lelbo^{\mathrm{ext}}_{V,\theta^{\mathrm{new}}}(\mathbb Q_{V,o,X};X,o)
\geq\Lelbo^{\mathrm{ext}}_{V,\theta^{\mathrm{old}}}(\mathbb Q_{V,o,X};X,o)
=\log p_{\theta^{\mathrm{old}},V}(o\given X).
\label{eq:elbo-monotonicity}
\end{equation}
Consequently a schedule that places a full exact E phase before each accepted
M phase has nondecreasing evidence values, which converge in
$(-\infty,+\infty]$ and to a finite limit when bounded above. Block
E-coordinates and accepted M-coordinates make the compared bound values
nondecreasing on a common functional; the transfer to the evidence in
\eqref{eq:elbo-monotonicity} uses the full E phase. \status{ESTABLISHED}

\noindent\emph{Proof.}
The first inequality is \cref{thm:elbo-extended-gap} at $\theta^{\mathrm{new}}$,
the second is \eqref{eq:elbo-acceptance}, and the equality is the zero-gap case
at $\theta^{\mathrm{old}}$; the singular-branch sum $\Lelbo^{\mathrm{ext}}+\KL$
is never written. A nondecreasing sequence of extended reals converges. Each
block coordinate does not decrease the bound by
\cref{prop:elbo-exact-e-coordinate}, and each accepted step does not decrease it
by \eqref{eq:elbo-acceptance}; only the final equality of
\eqref{eq:elbo-monotonicity} uses the exact E phase. Under (H3)--(H4) at both
candidates this is the classical conclusion of \citet{Dempster1977} in the
reading of \citet{Neal1998}. \hfill$\square$

Convergence of the parameter sequence needs regularity conditions of the kind
identified for the classical algorithm by \citet{Wu1983}. A finite step along a
natural-gradient or other optimizer direction is a proposal, not an accepted
coordinate, and it can decrease the bound.

\propositionheading{Finite natural-gradient steps need not be monotone}{prop:elbo-finite-step-nonmonotonicity}
Let $\mathsf G\succ0$, write $\nabla\Lelbo_V$ for the gradient of
$\theta'\mapsto\Lelbo_{V,\theta'}(\mathbb Q_{V,o,X};X,o)$, and consider the step
$\theta^+=\theta'+\mathrm h\mathsf G^{-1}\nabla\Lelbo_V$ of length
$\mathrm h>0$. At a nonstationary point the direction is an ascent direction,
\begin{equation}
\frac{d}{d\mathrm h}\Lelbo_{V,\theta^+}\Big|_{\mathrm h=0}
=\nabla\Lelbo_V^{\top}\mathsf G^{-1}\nabla\Lelbo_V>0.
\label{eq:elbo-ascent-direction}
\end{equation}
If $\Lelbo_{V,\theta''}=\Lelbo^\star-\tfrac12(\theta''-\theta^\star)^\top\mathsf H(\theta''-\theta^\star)$
with $\mathsf H\succ0$ for $\theta''$ in a neighborhood $\mathcal U$, $d_{\max}$
is the largest generalized eigenvalue of $(\mathsf H,\mathsf G)$, and
$\mathbf z$ is a corresponding eigenvector with
$\mathbf z^\top\mathsf G\mathbf z=1$, then at every
$\theta'=\theta^\star+\mathrm c_1\mathbf z\in\mathcal U$ with $\mathrm c_1\ne0$
and every step whose endpoint lies in $\mathcal U$,
\begin{equation}
\Lelbo_{V,\theta^+}-\Lelbo_{V,\theta'}
=\tfrac12d_{\max}\mathrm c_1^2\bigl[1-(1-\mathrm hd_{\max})^2\bigr]<0
\qquad\text{whenever}\qquad
\mathrm h>\frac{2}{d_{\max}}.
\label{eq:elbo-step-size-failure}
\end{equation}
\status{ESTABLISHED}

\noindent\emph{Proof.}
The first display is the chain rule. With
$\Delta\theta=\theta'-\theta^\star=\mathrm c_1\mathbf z$,
$\nabla\Lelbo_V=-\mathsf H\Delta\theta$ and
$\mathsf H\mathbf z=d_{\max}\mathsf G\mathbf z$ give
$\theta^+-\theta^\star=(1-\mathrm hd_{\max})\mathrm c_1\mathbf z$; substituting
in the quadratic gives the second display, whose bracket is negative exactly
when $\mathrm hd_{\max}>2$. \hfill$\square$

Monotonicity at a finite step therefore needs a line search, a trust region,
or another rule that verifies \eqref{eq:elbo-acceptance} on the same
functional; a better preconditioner does not remove the requirement, because
\eqref{eq:elbo-step-size-failure} is stated in the preconditioned spectrum.

\section{What remains open}
\label{sec:elbo-open}

\openproblemheading{Convergence of the alternating scheme}{open:elbo-alternating-convergence}
\Cref{prop:elbo-evidence-monotonicity} gives monotone convergence of the
evidence values under a full exact E phase and nothing about the parameters,
and since \eqref{eq:elbo-acceptance} is a weak inequality the null step is an
accepted M-coordinate, so the general scheme of block E-coordinates
\eqref{eq:elbo-e-coordinate} and accepted M-coordinates need not move its
parameter at all. The problem concerns models of
\cref{hyp:local-interaction-kernels} with finitely many latent states,
parameters in an open cone of a Euclidean space, and joint densities twice
continuously differentiable in the parameters, under the scheme that
alternates a full exact E phase with the exact M-coordinate
\eqref{eq:elbo-m-coordinate} whose maximizer is unique. The regularity
conditions of \citet{Wu1983} concern the limit points of the parameter
sequence and its convergence when the stationary points at its limit value
are isolated; isolation is not among the hypotheses of the class, in which
uniqueness of the maximizer leaves open that the evidence is constant along
a curve of parameters on which the complete-data likelihood is not. It is
settled by conditions checkable on the model, including compactness of the
upper level sets of the evidence in the cone, that are met by some model of
the class with parameter dimension at least two whose stationary set at the
limit value of an orbit is not discrete, under which every orbit of the
scheme converges to one stationary point of the extended bound, or by a
model of that class with compact upper level sets whose orbit has more than
one limit point.

\section*{Physicist's summary}

The population carries one generative law and one recognition law, and the
only evidence bound is $\log p(o)-\KL(Q\Vert\boldsymbol\Pi)$, defined for every
$Q$ and equal to $-\infty$ when $Q$ charges what the posterior cannot see;
everything else is bookkeeping of that relative entropy through the chain rule.
Substituting marginal entropies overstates the bound by exactly the total
correlation. Observations are records attached to latent states, and an
environment can be rewritten as boundary nodes emitting them. Each block of
agents has a local free energy built from the records that touch it; with the
rest held fixed, lowering it lowers the collective free energy by the same
amount, so the population is an exact potential game. Summed, the local free
energies count every shared record once per participant. Parameters move only
between evaluations, by accepted steps, and the evidence rises monotonically
when each step follows an exact posterior refresh; convergence is open.
